\documentclass[11pt, a4paper]{article}

% Gemeinsame Style-Definitionen (TEXINPUTS muss templates/ enthalten — siehe scripts/build_tex.py)
% bfw-style.tex — gemeinsame Style-Definitionen für alle Handouts/Aufgabenhefte
% Voraussetzung: Kompilation mit xelatex (wegen fontspec)

\usepackage[ngerman]{babel}
\usepackage{geometry}
\geometry{a4paper, top=2.2cm, bottom=2.2cm, left=2.3cm, right=2.3cm}

\usepackage{fontspec}
% Fallback-Kette: Source Sans Pro -> Calibri -> Segoe UI -> Latin Modern Sans
% (Latin Modern ist immer mit MiKTeX vorhanden)
\IfFontExistsTF{Source Sans Pro}{%
  \setmainfont{Source Sans Pro}[Ligatures=TeX]%
}{\IfFontExistsTF{Calibri}{%
  \setmainfont{Calibri}[Ligatures=TeX]%
}{\IfFontExistsTF{Segoe UI}{%
  \setmainfont{Segoe UI}[Ligatures=TeX]%
}{%
  \setmainfont{Latin Modern Sans}[Ligatures=TeX]%
}}}
\IfFontExistsTF{Source Code Pro}{%
  \setmonofont{Source Code Pro}[Scale=0.92]%
}{\IfFontExistsTF{Consolas}{%
  \setmonofont{Consolas}[Scale=0.92]%
}{%
  \setmonofont{Latin Modern Mono}[Scale=0.92]%
}}

\usepackage{xcolor}
% Farbpalette: hell, freundlich, modern
\definecolor{bfwPrimary}{HTML}{0EA5A4}    % Petrol/Türkis - Primär
\definecolor{bfwPrimaryDark}{HTML}{0F766E} % dunkler Petrol für Headings
\definecolor{bfwAccent}{HTML}{F59E0B}     % Amber - Akzent
\definecolor{bfwSuccess}{HTML}{10B981}    % Grün - Erfolg / Lösung
\definecolor{bfwWarn}{HTML}{F97316}       % Orange - Achtung
\definecolor{bfwInk}{HTML}{0F172A}        % fast Schwarz
\definecolor{bfwInkSoft}{HTML}{475569}    % gedämpftes Grau
\definecolor{bfwBgSoft}{HTML}{F8FAFC}     % off-white
\definecolor{bfwBgTeal}{HTML}{ECFEFF}     % helles Türkis
\definecolor{bfwBgAmber}{HTML}{FEF3C7}    % helles Gelb
\definecolor{bfwBgRose}{HTML}{FFE4E6}     % helles Rosé für Eselsbrücken

\usepackage[colorlinks=true, linkcolor=bfwPrimaryDark, urlcolor=bfwPrimaryDark]{hyperref}
\usepackage{enumitem}
\setlist[itemize]{leftmargin=*, itemsep=2pt, topsep=2pt}
\setlist[enumerate]{leftmargin=*, itemsep=2pt, topsep=2pt}

\usepackage[most]{tcolorbox}
\usepackage{booktabs}
\usepackage{tikz}
\usetikzlibrary{decorations.pathmorphing, positioning, shapes.geometric, arrows.meta}

% Merkkästchen — wichtige Definition / Kernpunkt
\newtcolorbox{merksatz}[1][]{%
  enhanced, breakable,
  colback=bfwBgTeal,
  colframe=bfwPrimary,
  boxrule=0.6pt,
  arc=4pt,
  left=10pt, right=10pt, top=8pt, bottom=8pt,
  fonttitle=\bfseries\color{bfwPrimaryDark},
  title={\faLightbulb\ Merksatz},
  #1
}

% Eselsbrücke — Mnemoniks
\newtcolorbox{eselsbruecke}[1][]{%
  enhanced, breakable,
  colback=bfwBgRose,
  colframe=bfwAccent,
  boxrule=0.6pt,
  arc=4pt,
  left=10pt, right=10pt, top=8pt, bottom=8pt,
  fonttitle=\bfseries\color{bfwWarn},
  title={Eselsbrücke},
  #1
}

% Beispiel-Box
\newtcolorbox{beispiel}[1][]{%
  enhanced, breakable,
  colback=bfwBgSoft,
  colframe=bfwInkSoft,
  boxrule=0.4pt,
  arc=3pt,
  left=10pt, right=10pt, top=8pt, bottom=8pt,
  fonttitle=\bfseries\color{bfwInk},
  title={Beispiel},
  #1
}

% Lösungs-Box (für Aufgabenhefte)
\newtcolorbox{loesung}[1][]{%
  enhanced, breakable,
  colback=bfwBgAmber!40,
  colframe=bfwSuccess,
  boxrule=0.6pt,
  arc=3pt,
  left=10pt, right=10pt, top=8pt, bottom=8pt,
  fonttitle=\bfseries\color{bfwSuccess!70!black},
  title={Lösung},
  #1
}

% Glossar-Eintrag
\newcommand{\glossareintrag}[2]{%
  \par\noindent\textbf{\color{bfwPrimaryDark}#1} \enspace #2\par\smallskip
}

% Section-Stil — etwas farbiger als Standard
\usepackage{titlesec}
\titleformat{\section}
  {\Large\bfseries\color{bfwPrimaryDark}}
  {\thesection}{0.6em}{}
\titleformat{\subsection}
  {\large\bfseries\color{bfwPrimary}}
  {\thesubsection}{0.5em}{}
\titleformat{\subsubsection}
  {\normalsize\bfseries\color{bfwInk}}
  {\thesubsubsection}{0.4em}{}

% TikZ-Stil "skizzenhaft" — etwas weicher als Rough.js, aber stabil
\tikzset{
  roughbox/.style = {
    draw=bfwPrimaryDark, thick, fill=bfwBgTeal,
    rounded corners=4pt, inner sep=6pt
  },
  roughaccent/.style = {
    draw=bfwAccent, thick, fill=bfwBgAmber,
    rounded corners=4pt, inner sep=6pt
  },
  roughline/.style = {
    draw=bfwInkSoft, thick, -{Stealth[length=2.5mm]}
  }
}

% Bullet-Symbole (bewusst kein FontAwesome — vermeidet Font-Installations-Probleme)
\newcommand{\faLightbulb}{$\bullet$}
\newcommand{\faBookmark}{$\bullet$}

% Header/Footer
\usepackage{fancyhdr}
\pagestyle{fancy}
\fancyhf{}
\renewcommand{\headrulewidth}{0pt}
\fancyfoot[L]{\small\color{bfwInkSoft}\thepage}
\fancyfoot[R]{\small\color{bfwInkSoft} BFW M-064 Beschaffung im Büromanagement}


\title{\vspace*{-1.5cm}\Huge\bfseries\color{bfwPrimaryDark}Tag~12: Dokumentenmanagement --- Aufbewahrung \& Fristen\\[0.4em]
  \large\color{bfwInkSoft}\textnormal{Das Gedächtnis des Unternehmens ordnen und sichern}}
\author{\textsf{Büroprozesse gestalten und Arbeitsvorgänge organisieren \textperiodcentered{} M-067}\\
\textsf{\small\copyright\ Can Siebert\,\textperiodcentered\,2026}}
\date{14.07.2026}

\begin{document}
\maketitle
\thispagestyle{empty}

\vspace{1em}
\begin{merksatz}[title={\bfseries Worum geht es heute?}]
Jeden Tag fallen in einem Unternehmen unzählige Schriftstücke an --- Rechnungen,
Verträge, Angebote, Kontoauszüge. Nicht alles darf weggeworfen, aber auch nicht alles
muss ewig aufbewahrt werden. In diesem Handout lernen Sie, \emph{warum} ein Betrieb
Schriftgut aufbewahrt, wie man Dokumente nach ihrer \emph{Wertigkeit} einstuft, welche
\emph{gesetzlichen Aufbewahrungsfristen} nach HGB und AO gelten --- und wo das Schriftgut
sinnvoll \emph{abgelegt} wird. Das ist die Grundlage für ein geordnetes
Dokumentenmanagement und ein wiederkehrendes Thema der IHK-Abschlussprüfung.
\end{merksatz}

\tableofcontents
\vspace{1em}
\hrule
\vspace{1em}

\section{Warum ein Betrieb Schriftgut aufbewahrt}

In einem Unternehmen liegt jeden Tag eine Vielzahl an Informationen und Daten in
geschriebener oder gedruckter Form vor --- auf Papier oder auf Datenträgern. Dazu
gehören etwa Eingangs- und Ausgangsrechnungen, Mahnungen, Arbeitsverträge,
Kündigungen, Anfragen, Angebote, Kaufverträge und Kontoauszüge.

Damit auf diese Unterlagen schnell und unkompliziert zurückgegriffen werden kann,
ergibt sich die \emph{innerbetriebliche Notwendigkeit} einer geordneten Aufbewahrung.
Darüber hinaus gibt es \emph{gesetzliche Vorschriften}, die die Aufbewahrung regeln: Die
Unterlagen dienen der \emph{Dokumentation und Beweissicherung} der Geschäftsvorfälle im
Rahmen der Buchhaltung. Zu diesem Zweck sollen die Buchführungsunterlagen innerhalb
bestimmter Fristen für Prüf- und Beweiszwecke verfügbar sein.

\begin{merksatz}
Man unterscheidet drei Gründe für die Aufbewahrung: \textbf{betriebliche} Gründe (schneller
Zugriff auf Vorgänge), \textbf{rechtliche} Gründe (Dokumentation und Beweis nach HGB) und
\textbf{steuerliche} Gründe (Prüfbarkeit durch das Finanzamt nach AO).
\end{merksatz}

Die Schriftgutverwaltung --- auch \emph{Registratur} genannt --- könnte man als das
\emph{„Gedächtnis eines Unternehmens”} bezeichnen. Damit dieses Gedächtnis am
Arbeitsplatz effektiv nutzbar ist, sollten bei der Verwaltung von Schriftgut u.\,a.\
folgende Überlegungen angestellt werden:

\begin{itemize}
  \item \textbf{Wertigkeit/Frist} --- Wie wichtig ist die Aufbewahrung; gibt es gesetzliche Vorgaben?
  \item \textbf{Ablageort} --- An welchem Ort im Unternehmen sollen die Schriftstücke aufbewahrt werden?
  \item \textbf{Ablagetechnik} --- Sollen die Schriftstücke lose oder geheftet abgelegt werden?
  \item \textbf{Aktenart} --- Sollen sie als Einzel- oder als Sammelakte abgelegt werden?
  \item \textbf{Ordnungssystem} --- Nach welchem logischen Prinzip wird abgeheftet?
\end{itemize}

\subsection{GoBD --- ordnungsmäßige (auch digitale) Aufbewahrung}

Werden Unterlagen elektronisch geführt oder gespeichert, müssen die \emph{GoBD} beachtet
werden: die \emph{Grundsätze zur ordnungsmäßigen Führung und Aufbewahrung von Büchern,
Aufzeichnungen und Unterlagen in elektronischer Form sowie zum Datenzugriff}. Nach
HGB § 257 Abs.~3 dürfen Unterlagen (außer Eröffnungsbilanzen und Abschlüssen) auch als
Wiedergabe auf einem Bild- oder Datenträger aufbewahrt werden, wenn dies den Grundsätzen
ordnungsmäßiger Buchführung entspricht und sichergestellt ist, dass die Wiedergabe

\begin{itemize}
  \item mit den empfangenen Handelsbriefen und Buchungsbelegen \emph{bildlich} und mit den
    übrigen Unterlagen \emph{inhaltlich} übereinstimmt, wenn sie lesbar gemacht wird, und
  \item während der Aufbewahrungsfrist jederzeit innerhalb angemessener Frist verfügbar
    und lesbar gemacht werden kann.
\end{itemize}

Zentral ist die \emph{Unveränderbarkeit}: Nur ein nachträglich nicht manipulierbarer
Beleg hat Beweiskraft und wird vom Finanzamt anerkannt.

\section{Wertigkeitsstufen von Schriftgut}

Die anfallenden Schriftstücke lassen sich je nach Bedeutung für den Betrieb
unterschiedlichen \emph{Wertigkeitsstufen} zuordnen. In der Praxis unterscheidet man
vier Stufen:

\begin{enumerate}
  \item \textbf{Tageswert} --- kein bleibender Wert für das Unternehmen. Beispiele:
    Tageszeitungen, Prospekte, unverlangte Angebote. Sie werden nach Kenntnisnahme
    vernichtet (interessante Artikel können ausgeschnitten werden).
  \item \textbf{Prüfwert} --- nur für einen \emph{begrenzten Zeitraum} von Bedeutung. Nach
    der Prüfung wird entschieden, ob vernichtet oder aufbewahrt wird. Beispiele:
    verlangte Angebote, Preislisten, Statistiken.
  \item \textbf{Gesetzeswert} --- der Gesetzgeber schreibt im HGB und in der AO
    \emph{Aufbewahrungsfristen} vor, damit ein Handelsgeschäft lückenlos nachvollzogen
    werden kann. Beispiele: Rechnungen, Buchungsbelege, Jahresabschlüsse.
  \item \textbf{Dauerwert} --- von besonderer Bedeutung; aufzubewahren, solange das
    Unternehmen besteht. Beispiele: Unterlagen zur Unternehmensgründung, Patentunterlagen.
\end{enumerate}

\begin{eselsbruecke}
\textbf{„Tagsüber Prüfe Gesetze Dauerhaft”} --- die Anfangsbuchstaben \emph{T-P-G-D}
erinnern an \emph{Tages-, Prüf-, Gesetzes-} und \emph{Dauerwert} --- vom kürzesten bis
zum längsten Aufbewahrungsbedarf.
\end{eselsbruecke}

\section{Aufbewahrungsfristen nach HGB und AO}

Schriftstücke mit \emph{Gesetzeswert} müssen für eine gesetzlich festgelegte Dauer
aufbewahrt werden. Die zentrale Vorschrift ist \emph{HGB § 257 (Aufbewahrung von
Unterlagen. Aufbewahrungsfristen)}. Für das Steuerrecht enthält \emph{AO § 147}
entsprechende, weitgehend deckungsgleiche Regelungen --- diese erfassen auch
Unternehmen, die keine Kaufleute nach HGB sind, aber zur Buchführung verpflichtet sind
oder freiwillig Bücher führen.

\subsection{Fristen und Fristbeginn}

\begin{center}
\begin{tabular}{p{8.6cm} c}
\toprule
\textbf{Art der Unterlage} & \textbf{Frist}\\
\midrule
Handelsbücher, Inventare, Eröffnungsbilanzen, Jahresabschlüsse, Lageberichte & 10 Jahre\\
Buchungsbelege (z.\,B.\ Rechnungen, Quittungen, Kontoauszüge) & siehe Hinweis\\
Empfangene und abgesandte Handelsbriefe & 6 Jahre\\
Sonstige Unterlagen mit steuerlicher Bedeutung & 6 Jahre\\
\bottomrule
\end{tabular}
\end{center}

\medskip
Ein \emph{Handelsbrief} ist dabei jedes Schriftstück, das ein Handelsgeschäft betrifft.
Nicht alle Geschäftsbriefe zählen dazu --- z.\,B.\ unverbindliche Werbung nicht.

\begin{merksatz}[title={\bfseries Wichtig: Fristbeginn}]
Die Aufbewahrungsfrist beginnt mit dem \emph{Schluss des Kalenderjahres}, in dem die
letzte Eintragung gemacht, der Jahresabschluss festgestellt, der Handelsbrief empfangen
oder abgesandt oder der Buchungsbeleg entstanden ist. Das Fristende liegt deshalb immer
auf einem \textbf{31.\ Dezember}.
\end{merksatz}

\subsection{Aktueller Hinweis: Verkürzung auf 8 Jahre (BEG IV)}

Durch das \emph{Vierte Bürokratieentlastungsgesetz (BEG~IV)}, das seit 2025 in Kraft ist,
wurde die Aufbewahrungsfrist für \emph{Buchungsbelege} (HGB § 257 Abs.~4, AO § 147 Abs.~3)
von zehn auf \textbf{acht Jahre} verkürzt. Wichtig zur Einordnung:

\begin{itemize}
  \item \textbf{Buchungsbelege} (z.\,B.\ Rechnungen, Quittungen, Kontoauszüge): \textbf{8 Jahre}.
  \item \textbf{Handelsbücher, Inventare, Eröffnungsbilanzen, Jahresabschlüsse, Lageberichte}:
    weiterhin \textbf{10 Jahre}.
  \item \textbf{Handelsbriefe} und sonstige steuerlich bedeutsame Unterlagen: weiterhin
    \textbf{6 Jahre}.
\end{itemize}

\begin{beispiel}
Eine Eingangsrechnung (Buchungsbeleg) entsteht am 12.\,03.\,2025. Die Frist beginnt mit
dem Schluss des Kalenderjahres, also am 31.\,12.\,2025. Bei einer Frist von 8 Jahren endet
die Aufbewahrungspflicht am \textbf{31.\,12.\,2033}. Ein im selben Jahr 2025 empfangener
Handelsbrief (6 Jahre) wäre dagegen bis zum \textbf{31.\,12.\,2031} aufzubewahren.
\end{beispiel}

\subsection{Der Schriftgutkatalog}

Zur Vermeidung von Doppel- und Mehrfachablagen sowie eines unnötig langen Aufbewahrens
sollte in jedem Büro ein \emph{Schriftgutkatalog} erstellt werden. Er verzeichnet alle
in diesem Büro vorkommenden Schriftstücke mit der jeweiligen \emph{Wertigkeitsstufe},
der \emph{Aufbewahrungsfrist} und dem \emph{Ablageort}.

\begin{center}
\begin{tabular}{p{3.4cm} p{3.0cm} p{2.0cm} p{3.0cm}}
\toprule
\textbf{Belegart} & \textbf{Wertigkeit} & \textbf{Frist} & \textbf{Ende am}\\
\midrule
Ausgangsrechnungen & Gesetzeswert & 8 Jahre & 31.12.2033\\
Jahresabschluss & Gesetzeswert & 10 Jahre & 31.12.2036\\
Angebote (verlangt) & Prüfwert & --- & Gültigkeit beachten\\
Preislisten/Kataloge & Prüfwert & --- & Gültigkeit beachten\\
Gründungsurkunde & Dauerwert & dauerhaft & ---\\
\bottomrule
\end{tabular}
\end{center}

\section{Ablageorte und Registraturarten}

Bei der Entscheidung, an welchem \emph{Ort} die Unterlagen aufbewahrt werden, gilt der
Grundsatz: \emph{Je häufiger ein Mitarbeiter auf die Schriftstücke zurückgreifen muss,
umso näher an seinem Arbeitsplatz sollten sie aufbewahrt werden.}

\subsection{Mögliche Ablageorte}

\begin{enumerate}
  \item \textbf{Arbeitsplatzablage} --- Schriftstücke, die ein Mitarbeiter aktuell
    bearbeitet oder auf die er häufig und schnell zugreifen muss, bewahrt er direkt am
    Arbeitsplatz auf (dezentral).
  \item \textbf{Abteilungsablage} --- Unterlagen, die von mehreren Mitarbeitern
    gleichzeitig bearbeitet oder verwaltet werden, lagern gut erreichbar in der
    Abteilung. Die Ablage muss für alle nachvollziehbar sein; hilfreich ist ein
    \emph{Aktenplan}.
  \item \textbf{Zentralregistratur/Archiv} --- hier werden Unterlagen verschiedener
    Bereiche aufbewahrt. Die Zentralregistratur entlastet die Abteilungen von
    Schriftstücken, die nicht mehr laufend benötigt werden, aber noch aufzubewahren sind
    (z.\,B.\ wegen der gesetzlichen Aufbewahrungspflicht). Auch Unterlagen mit Dauerwert
    werden hier archiviert.
\end{enumerate}

\subsection{Registraturarten: zentral und dezentral}

Nach dem Aktualitätsgrad der Schriftstücke unterscheidet man drei Registraturarten:

\begin{itemize}
  \item \textbf{Tagesregistratur} (laufende/aktive Registratur) --- Unterlagen aus dem
    aktuellen Geschäftsjahr, auf die laufend zugegriffen wird; meist dezentral am
    Arbeitsplatz oder in der Abteilung.
  \item \textbf{Zwischenregistratur} --- abgeschlossene Vorgänge, die noch gelegentlich
    gebraucht werden; ausgelagert, aber noch gut erreichbar.
  \item \textbf{Altregistratur} (Archiv) --- selten benötigte, aber noch
    aufbewahrungspflichtige Unterlagen sowie Dauerwert-Schriftgut; meist zentral.
\end{itemize}

Eine \emph{dezentrale} Ablage (am Arbeitsplatz/in der Abteilung) bietet schnellen
Zugriff, eine \emph{zentrale} Ablage (Zentralregistratur) sorgt für einheitliche
Ordnung und entlastet die Arbeitsplätze.

\section{Zusammenfassung}

\begin{enumerate}
  \item Betriebe bewahren Schriftgut aus \textbf{betrieblichen, rechtlichen und steuerlichen} Gründen auf; die Registratur ist das „Gedächtnis des Unternehmens”.
  \item Bei digitaler Aufbewahrung gelten die \textbf{GoBD} --- entscheidend sind Lesbarkeit, Verfügbarkeit und \textbf{Unveränderbarkeit}.
  \item Vier \textbf{Wertigkeitsstufen}: Tages-, Prüf-, Gesetzes- und Dauerwert.
  \item \textbf{Fristen} nach HGB § 257 / AO § 147: 10 Jahre (Bücher, Abschlüsse), 6 Jahre (Handelsbriefe); Buchungsbelege seit BEG~IV \textbf{8 Jahre}.
  \item Der \textbf{Fristbeginn} ist stets der Schluss des Kalenderjahres --- das Fristende immer ein 31.\,Dezember.
  \item \textbf{Ablageorte} richten sich nach der Zugriffshäufigkeit: Arbeitsplatz-, Abteilungsablage, Zentralregistratur/Archiv (Tages-, Zwischen-, Altregistratur).
\end{enumerate}

\newpage

\section*{Glossar}
\addcontentsline{toc}{section}{Glossar}

\glossareintrag{Schriftgut}{Alle in einem Unternehmen anfallenden Informationen und Daten in geschriebener oder gedruckter Form auf Papier oder Datenträgern.}

\glossareintrag{Schriftgutverwaltung (Registratur)}{Geordnete Aufbewahrung und Verwaltung des Schriftguts; das „Gedächtnis des Unternehmens”.}

\glossareintrag{GoBD}{Grundsätze zur ordnungsmäßigen Führung und Aufbewahrung von Büchern, Aufzeichnungen und Unterlagen in elektronischer Form sowie zum Datenzugriff.}

\glossareintrag{Tageswert}{Wertigkeitsstufe ohne bleibenden Wert; nach Kenntnisnahme vernichtbar (z.\,B.\ Tageszeitung, Prospekt).}

\glossareintrag{Prüfwert}{Wertigkeitsstufe von nur begrenzter Bedeutung; nach Prüfung wird über Vernichtung/Aufbewahrung entschieden (z.\,B.\ Preisliste).}

\glossareintrag{Gesetzeswert}{Wertigkeitsstufe mit gesetzlicher Aufbewahrungsfrist nach HGB/AO (z.\,B.\ Rechnungen, Jahresabschlüsse).}

\glossareintrag{Dauerwert}{Wertigkeitsstufe von besonderer Bedeutung; aufzubewahren, solange das Unternehmen besteht (z.\,B.\ Gründungsunterlagen).}

\glossareintrag{HGB § 257}{Vorschrift des Handelsgesetzbuchs zur Aufbewahrung von Unterlagen und zu den Aufbewahrungsfristen.}

\glossareintrag{AO § 147}{Vorschrift der Abgabenordnung mit weitgehend deckungsgleichen steuerlichen Aufbewahrungsregelungen.}

\glossareintrag{Handelsbrief}{Jedes Schriftstück, das ein Handelsgeschäft betrifft; Aufbewahrungsfrist 6 Jahre.}

\glossareintrag{Buchungsbeleg}{Beleg für eine Buchung (z.\,B.\ Rechnung, Quittung, Kontoauszug); Aufbewahrungsfrist seit BEG~IV 8 Jahre.}

\glossareintrag{Fristbeginn}{Schluss des Kalenderjahres, in dem das Schriftstück entstanden bzw.\ empfangen/abgesandt wurde; Fristende stets 31.\,Dezember.}

\glossareintrag{BEG IV}{Viertes Bürokratieentlastungsgesetz; verkürzte seit 2025 die Aufbewahrungsfrist für Buchungsbelege von 10 auf 8 Jahre.}

\glossareintrag{Schriftgutkatalog}{Verzeichnis aller Schriftstücke eines Büros mit Wertigkeitsstufe, Aufbewahrungsfrist und Ablageort; vermeidet Doppelablagen.}

\glossareintrag{Arbeitsplatzablage}{Dezentrale Ablage direkt am Arbeitsplatz für häufig und schnell benötigte Unterlagen.}

\glossareintrag{Abteilungsablage}{Gemeinsam genutzte Ablage einer Abteilung; für alle nachvollziehbar (Aktenplan).}

\glossareintrag{Zentralregistratur/Archiv}{Zentrale Aufbewahrung von Unterlagen verschiedener Bereiche; entlastet Abteilungen und bewahrt Dauerwert-Schriftgut.}

\glossareintrag{Aktenplan}{Systematische Gliederung, die eine übersichtliche und nachvollziehbare Ordnung des Schriftguts sicherstellt.}

\end{document}
